\chapter{Gerçekleştirme}
\label{ch:gerceklestirme}

Bu bölümde, \cref{ch:yontem}'de tanımlanan tasarımın somut donanım ve
yazılım bileşenlerine indirgenişi anlatılmaktadır. Her alt sistem için
\emph{ne kullanıldı, neden kullanıldı, nasıl bağlandı, hangi
parametrelerle çalıştığı} sırasıyla verilmiştir.

\section{Donanım Listesi}

\Cref{tab:bom} prototipte kullanılan tüm donanım bileşenlerini ve
yaklaşık birim maliyetlerini özetler. Toplam ek donanım maliyeti,
hedeflenen 500~TL bütçesi içinde kalmıştır.

\begin{table}[h!]
\centering
\caption{Bileşen listesi (BOM) ve yaklaşık birim maliyetler.}
\label{tab:bom}
\small
\begin{tabular}{l l c l}
\toprule
\textbf{Parça} & \textbf{Model / Tip} & \textbf{Adet} & \textbf{Tahmini Fiyat}\\
\midrule
Olay yöneticisi MCU      & STM32 NUCLEO-F767ZI      & 1 & (Mevcut, bütçeye eklenmedi)\\
Kamera + Wi-Fi modülü    & ESP32-CAM AI-Thinker (OV3660) & 1 & 120 TL\\
BLE köprü modülü         & HM-10 (CC2541)           & 1 & 80 TL\\
LED                      & Yeşil 5\,mm, sarı, kırmızı + 220\,$\Omega$ direnç & 3+3 & 10 TL\\
Buzzer                   & Aktif buzzer 5\,V        & 1 & 15 TL\\
Buton                    & Push button taktil 6\,mm & 2 & 10 TL\\
Breadboard + jumper      & --                       & -- & 100 TL\\
USB-TTL dönüştürücü      & CP2102 / FT232           & 1 & 50 TL\\
Yedek 40-pin male header & --                       & 2 & 10 TL\\
\midrule
\textbf{Toplam ek donanım} & & & \textbf{$\approx$ 395 TL}\\
\bottomrule
\end{tabular}
\end{table}

Buluttaki çalışma maliyeti için AWS EC2 \texttt{m7i-flex.large}
örneği eu-central-1 (Frankfurt) bölgesinde çalıştırılmaktadır;
AWS Başlangıç Kredisi ile sergi ve test dönemi için ek maliyet
oluşturmamıştır.

\section{Sunucu Tarafı: Flask + InsightFace}
\label{sec:sunucu-impl}

\subsection{Yığın ve Bağımlılıklar}

Sunucu tarafı Python 3.11 ile yazılmış, Flask web framework'ü ve
InsightFace 0.7 paketi üzerine kurulmuştur. ONNX Runtime CPU sürümü
kullanılmaktadır; bu sayede üretim ortamı (Linux/EC2) ve geliştirme
ortamı (macOS) arasında deterministik davranış sağlanmaktadır.
Konteynerleştirme Docker ile yapılmış; Gunicorn $-w\,1$ yapılandırmasıyla
çalıştırılmaktadır. Tek worker zorunluluğu, oturum (session) verisinin
süreç içi (in-memory) saklanmasından kaynaklanır: birden fazla worker,
aynı oturumun karelerini farklı süreçlere dağıtarak agregasyonu kırardı.

\subsection{Endpoint Tasarımı}

Tek bir kritik endpoint mevcuttur: \texttt{POST /search}. İstek
gövdesi JPEG bayt dizisi; başlıkları:

\begin{itemize}[itemsep=1pt]
    \item \texttt{X-Session-Id}: Oturum kimliği (UUID), tüm burst boyunca aynı.
    \item \texttt{X-Frame-Index}: 1\textendash N arası kare sırası.
    \item \texttt{X-Frame-Total}: Burst içindeki toplam kare sayısı (\=10).
    \item \texttt{X-Shared-Secret}: (Opsiyonel) İlk faz için kapalı bir
          anahtar; üretimde set edildiğinde \texttt{X-Shared-Secret}'i
          taşımayan istekler 401 ile reddedilir.
\end{itemize}

Ara karelere yanıt: \texttt{\{"status":"continue", ...\}}; son kareye
yanıt:

\begin{lstlisting}[language=none, caption={Son kareye sunucu yanıt örneği.}, label=lst:final-response]
{
  "match": true,
  "name": "Ali_Yilmaz",
  "similarity": 0.9437,
  "frames_used": 10,
  "frames_total": 10,
  "liveness_score": 0.014823
}
\end{lstlisting}

\subsection{Operasyonel Sertleştirmeler}

Sergi öncesi yapılan üç hafif sertleştirme:

\begin{enumerate}[itemsep=2pt]
    \item \textbf{Opsiyonel paylaşılan anahtar:} \texttt{SHARED\_SECRET}
          ortam değişkeni atandığında \texttt{/search} endpoint'i
          \texttt{X-Shared-Secret} başlığını arar; \texttt{/health} ve
          \texttt{/metrics} izleme için açık tutulur.
    \item \textbf{Yapılandırılmış JSON günlüğü:} Her burst için bir
          satırlık JSON kayıt (\texttt{event:"burst\_complete"}, oturum
          kimliği, match kararı, similarity, liveness, gecikme) standart
          çıktıya yazılır. Bu format, sergi sonrası analiz ve tez Bölüm 5
          tablolarının üretimi için doğrudan kullanılabilir.
    \item \textbf{\texttt{/metrics} endpoint'i:} Açık oturum sayısı,
          tamamlanan burst sayısı, MATCH/NOMATCH dağılımı, son 100
          burst'ün gecikme p50/p95/maks değerleri, uptime ve kimlik
          doğrulama durumu tek bir JSON cevabıyla raporlanır.
\end{enumerate}

\section{ESP32-CAM: Kamera ve Ağ Köprüsü}

\subsection{Donanım Yapılandırması}

AI-Thinker ESP32-CAM modülü, ESP32-WROOM-32 SoC üzerine entegre OV3660
kamera modülü ve antenden oluşur. Wi-Fi STA modunda sürücünün telefon
hotspot'una bağlanır; flash etmek için 3.3~V FTDI bağlantısı kullanılır
(IO0 GND'ye çekilerek bootloader moduna geçilir).

\subsection{Yazılım Akışı}

Firmware PlatformIO + Arduino-ESP32 çekirdeği üzerinde C++ olarak
yazılmıştır. Ana döngü iki görevi yürütür:

\begin{enumerate}[label=(\arabic*),itemsep=1pt]
    \item STM32'den UART üzerinden gelen \texttt{CAPTURE\textbackslash n}
          komutunu dinlemek,
    \item Komut geldiğinde 5~FPS hızında 10 JPEG kare yakalamak, her
          birini yeni bir HTTP POST isteğiyle EC2'ye iletmek ve son
          karenin JSON cevabını parse etmek.
\end{enumerate}

Burst boyunca her POST aynı oturum kimliğini (\texttt{X-Session-Id})
taşır. Son kare üzerindeki yanıt JSON olarak parse edildikten sonra:

\begin{itemize}[itemsep=1pt]
    \item Eşleşme varsa: STM32'ye \texttt{1;<name>;<sim>\textbackslash n} CSV satırı,
    \item Eşleşme yoksa: \texttt{0;;<sim>\textbackslash n},
    \item Hata durumunda: \texttt{ERR:<tag>\textbackslash n} (örnek tag:
          \texttt{no\_wifi}, \texttt{http\_500}, \texttt{json\_parse}).
\end{itemize}

\subsection{Yapılandırma}

Wi-Fi bilgileri, sunucu URL'si ve TLS bayrağı \texttt{include/config.h}
dosyasından okunur. Faz 1'de (HTTP) ve Faz 2'de (HTTPS + Let's Encrypt)
arasındaki tek farklılık bu dosyadır; sertifika güncellemesi ve TLS
bayrağı ayarı sonrasında yeni image flash edilir.

\section{STM32 Olay Yöneticisi: Gerçekleme}

\subsection{Donanım Konfigürasyonu}

CubeMX projesi NUCLEO-F767ZI üzerinde aşağıdaki çevre birimleri ile
yapılandırılmıştır:

\begin{itemize}[itemsep=1pt]
    \item UART1, 115200~bps 8N1, PA9/PA10: ESP32-CAM ile haberleşme,
    \item UART2, 9600~bps 8N1, PD5/PD6: HM-10 ile haberleşme,
    \item GPIO: PB0 (yeşil), PD12 (sarı), PB14 (kırmızı), PD13 (buzzer),
    \item EXTI: PC13 (TARA butonu, kart üstü mavi USER butonu), PA0 (harici PANİK butonu).
\end{itemize}

\textbf{Önemli not:} F767ZI'de Ethernet çevre birimi varsayılan olarak
aktif gelir ve DCMI pinleri ile çakışır. Bu projede DCMI kullanılmadığı
için ETH disabled bırakılmıştır; eski Plan A taslağındaki kamera doğrudan
DCMI'ye bağlanacak senaryosu Plan B'ye geçişle birlikte düşmüştür.

\subsection{Yazılım Mimarisi}

STM32 firmware'i iki katmanlı geliştirilmiştir.

\paragraph{(i) Tasarım doğrulama katmanı.}
Durum makinesi mantığı, taşınabilir C99 olarak yazılmış ayrı bir
\textbf{referans modül} (\texttt{state\_machine.c/.h}) hâlinde ifade
edilmiş ve PC üzerinde birim testleriyle doğrulanmıştır. Modül, HAL'a
doğrudan bağımlı değildir; tüm donanım etkileşimleri bir callback
vtable üzerinden gerçekleşir:

\begin{lstlisting}[language=C, caption={State machine referans modülünün callback vtable'ı.}, label=lst:cb]
typedef struct {
    void     (*set_led)(sm_led_t led, bool on);
    void     (*set_buzzer)(bool on);
    void     (*send_to_esp)(const char *line);
    void     (*send_to_hm10)(const char *line);
    uint32_t (*now_ms)(void);
} sm_callbacks_t;
\end{lstlisting}

Bu soyutlama sayesinde modülün davranışı donanım olmadan
gcc/clang ile derlenen bir host-side test koşumunda deterministik olarak
doğrulanmıştır. \Cref{tab:host-tests} kapsanan 12 senaryoyu özetler.

\paragraph{(ii) Embedded gerçekleme.}
Tasarım doğrulandıktan sonra, \texttt{main.c} dosyasında aynı durum
makinesi mantığı CubeMX'in yeniden ürettiği \emph{USER CODE} bölgelerine
\emph{inline} edilmiştir. Bu pratik tercih iki nedenden ötürü
yapılmıştır: (i) CubeMX yenileme süreci ile çakışma riskini azaltmak,
(ii) embedded tarafta UART RX kesme rutininin ve EXTI callback'lerinin
state machine eventleri ile aynı dosyada yaşamasının okunabilirliği
artırması. Embedded ve referans uygulamalar arasındaki davranış birebir
aynıdır; tek farklılık hold timer sabitleri (NOMATCH 1\,s vs.\ 2\,s
gibi), proje yürüyüşü sırasında pratiğe göre kalibre edilmiştir.

Bu iki katmanlı yaklaşım, donanım olmadan tasarım hatalarını yakalama
ve donanım üzerinde flash--flash döngülerini en aza indirme açısından
proje takvimini önemli ölçüde rahatlatmıştır.

\begin{table}[h!]
\centering
\caption{STM32 durum makinesi için host-side birim test kapsamı.}
\label{tab:host-tests}
\small
\begin{tabular}{l l}
\toprule
\textbf{Test senaryosu} & \textbf{Doğrulanan davranış}\\
\midrule
initial state              & IDLE + yeşil LED sabit\\
TARA press                 & SCANNING geçişi + CAPTURE komutu\\
ESP MATCH                  & MATCH durumu + kırmızı LED + buzzer + HM-10 frame\\
ESP NOMATCH                & NOMATCH yolu, yeşil LED, doğru frame\\
scan timeout               & 15~s sonrası NETERR\\
PANIC during scan          & PANIC durumunun her durumdan erişilebilir olması\\
hold expiry                & MATCH/NOMATCH/NETERR sonunda IDLE'a dönüş\\
heartbeat                  & IDLE'da 5~s aralıkla HB satırı\\
TARA during hold           & Hold penceresinin kesilebilmesi\\
late ESP result            & Timeout sonrası gelen yanıtın yok sayılması\\
NETERR blink               & NETERR'de sarı LED'in periyodik geçişi\\
oversized name             & 32 karakterden uzun isimde taşma yok\\
\bottomrule
\end{tabular}
\end{table}

\subsection{Embedded Uygulama --- main.c}

\texttt{main.c} dosyasında durum makinesi mantığı USER CODE bölgelerine
inline yazılmıştır. \Cref{lst:adapter} temel parçaları gösterir.

\begin{lstlisting}[language=C, caption={CubeIDE main.c içerisindeki durum makinesi gerçeklemesi (özet).}, label=lst:adapter]
typedef enum {
    STATE_IDLE = 0, STATE_SCANNING, STATE_MATCH,
    STATE_NOMATCH,  STATE_PANIC,    STATE_NETERR,
} app_state_t;

static volatile app_state_t app_state = STATE_IDLE;
static uint32_t state_entered_ms = 0;

static void apply_state_outputs(app_state_t s) {
    /* tüm LED + buzzer kapat, sonra duruma göre aç */
    HAL_GPIO_WritePin(LED_GREEN_GPIO_Port,  LED_GREEN_Pin,  GPIO_PIN_RESET);
    /* ... */
    switch (s) {
        case STATE_IDLE: HAL_GPIO_WritePin(LED_GREEN_GPIO_Port, LED_GREEN_Pin, GPIO_PIN_SET); break;
        case STATE_MATCH: case STATE_PANIC:
            HAL_GPIO_WritePin(LED_RED_GPIO_Port, LED_RED_Pin, GPIO_PIN_SET);
            HAL_GPIO_WritePin(BUZZER_GPIO_Port,  BUZZER_Pin,  GPIO_PIN_SET);
            break;
        /* ... */
    }
}

void HAL_GPIO_EXTI_Callback(uint16_t pin) {
    if (pin == BTN_SCAN_Pin)  flag_tara  = true;
    if (pin == BTN_PANIC_Pin) flag_panic = true;
}
\end{lstlisting}

UART1 RX yolu (ESP'den gelen \texttt{RESULT:..}/\texttt{ERR:..}/\texttt{HB}
satırları), bir bayt-bayt kesme rutini ile \texttt{\textbackslash n}'a kadar bir
satır tamponuna birikir; ana döngüde tamamlanan satırlar
\texttt{process\_cam\_line()} ile parse edilerek MATCH/NOMATCH/NETERR
geçişlerini tetikler.

\section{Android Uygulaması: TaxiGuard v2}

\subsection{Bileşenler}

Uygulama Kotlin + Jetpack Compose ile yazılmış, hedef SDK 34, minimum
SDK 26'dır. Üç ana modülden oluşur:

\begin{itemize}[itemsep=2pt]
    \item \textbf{\texttt{Hm10Service}}: Bir foreground servis olarak
          BLE central rolünde çalışır. HM-10 modülünü \texttt{0xFFE0}
          servis UUID'sine göre tarar, bulunduğunda GATT bağlantısı kurar
          ve \texttt{0xFFE1} karakteristiğinin notifikasyonlarına abone
          olur. MTU parçalanmasına karşı dahili bir \emph{line buffer}
          her \texttt{\textbackslash n}'a kadar baytları biriktirir.
    \item \textbf{\texttt{FrameParser}}: Gelen ASCII satırları parse
          ederek bir \texttt{MatchEvent} sealed hierarchy'sine dönüştürür
          (\texttt{Scanning, Match(name, sim), NoMatch, Panic, NetError,
          Heartbeat, Unknown}).
    \item \textbf{\texttt{EmergencyDialer}}: MATCH veya PANIC event'i
          alındığında \texttt{Intent.ACTION\_CALL("tel:155")}'i çağırır;
          \texttt{CALL\_PHONE} izni yoksa \texttt{ACTION\_DIAL} dönüşüne
          fallback yapar (dialer önceden doldurulmuş şekilde açılır).
\end{itemize}

UI tek ekrandır (\texttt{HomeScreen}); bağlantı bandı, son MATCH kartı,
ekran içi olay günlüğü ve servis Başlat/Durdur kontrolleri içerir.
Ayrıca, donanım yokken bile parse + dial akışını tek telefonda doğrulamak
için \emph{``DEV: Simulate MATCH''} butonu yer alır. Bu buton, servise
\texttt{ACTION\_SIMULATE\_MATCH} intent'i göndererek sahte bir
\texttt{MATCH:Test\_Subject;0.95} satırını parse pipeline'ına besler;
böylece HM-10 elimizde olmadan da uygulamanın MATCH$\rightarrow$155
zinciri doğrulanabilmektedir.

\subsection{İzinler ve Sandbox Davranışı}

\Cref{tab:android-izinler} uygulamanın manifest'inde tanımlı izinleri ve
gerekçelerini özetler.

\begin{table}[h!]
\centering
\caption{Android izinleri ve gerekçeleri.}
\label{tab:android-izinler}
\small
\begin{tabular}{l l}
\toprule
\textbf{İzin} & \textbf{Gerekçe}\\
\midrule
\texttt{BLUETOOTH\_SCAN}                          & API 31+ BLE cihaz keşfi\\
\texttt{BLUETOOTH\_CONNECT}                        & GATT bağlantı kurma\\
\texttt{CALL\_PHONE}                               & \texttt{Intent.ACTION\_CALL("tel:155")}\\
\texttt{FOREGROUND\_SERVICE\_CONNECTED\_DEVICE}    & Servis kategorisi (BLE)\\
\texttt{POST\_NOTIFICATIONS}                       & API 33+ foreground bildirimi\\
\texttt{ACCESS\_FINE\_LOCATION}                    & API 26--30 (legacy BLE keşif)\\
\bottomrule
\end{tabular}
\end{table}

\section{Ölçüm Harness'i: \texttt{far\_frr.py}}
\label{sec:harness}

Tez \cref{ch:sonuclar} ölçümlerinin tekrarlanabilir biçimde
yapılabilmesi için \texttt{eval/far\_frr.py} adlı bir Python aracı
geliştirilmiştir. Araç bir veri seti klasörünü (\texttt{dataset/<isim>/<foto>.jpg})
tarar, her fotoğrafı 10 karelik bir burst olarak \texttt{/search}
endpoint'ine gönderir ve dönen similarity skorlarını CSV'ye yazar.
Ardından eşik $\tau \in [0, 1]$ üzerinde 0.01 adımlarla bir süpürme
gerçekleştirerek FAR, FRR ve TA/FA/TR/FR sayılarını üretir; Equal Error
Rate (EER) eşiğini ve istenen maksimum FAR altındaki en iyi operating
point'i raporlar. Matplotlib mevcutsa bir FAR--FRR grafiği de
oluşturur.

\Cref{tab:harness-cikti} aracın ürettiği çıktı dosyalarını özetler.

\begin{table}[h!]
\centering
\caption{\texttt{far\_frr.py} aracının ürettiği çıktı dosyaları.}
\label{tab:harness-cikti}
\small
\begin{tabular}{l l}
\toprule
\textbf{Dosya} & \textbf{İçerik}\\
\midrule
\texttt{<run>.csv}            & Her probe için bir satır; similarity, matched name, kategori\\
\texttt{<run>\_sweep.csv}     & Eşik--FAR/FRR tablosu (101 satır)\\
\texttt{<run>\_summary.json}  & EER eşiği, operating point, sayımlar\\
\texttt{<run>\_roc.png}       & FAR--FRR eğrisi (matplotlib varsa)\\
\bottomrule
\end{tabular}
\end{table}

Aracın ürettiği CSV ve JSON dosyaları, tezin tablo ve grafiklerine
doğrudan kaynak olur; \texttt{run\_yyyy\_mm\_dd} isimlendirmesi sayesinde
hangi koşulda yapılmış ölçüm olduğu izlenebilir.
